\section{Background and Related Work}
\label{sec:background}

% PAPER-FLOW §2.1–§2.3
% Covers: (a) IP geolocation toolkit comparison (DNS, Geofeed, traceroute, latency/CBG),
% emphasizing that latency-based methods double as validation for the other three;
% (b) CBG at a glance — three-phase pipeline (LTD → MTL → CTR), inputs/outputs;
% (c) the four foundational variants in a table: Vanilla (low-envelope + planar-circle +
% geometric-centroid), Million-scale (speed-of-internet + planar-circle + geometric-centroid),
% Octant (bounded-spline + planar-annulus-weighted + Monte-Carlo), Spotter (normal-dist +
% planar-annulus-weighted + Monte-Carlo); implementation note that Vanilla/Million-scale
% use PlanarCircleMTL + geometric_centroid (not spherical-circle + BVM as in the originals).
%
% Key citations: Katz-Bassett 2006 (CBG), Laki 2011 (Spotter), Poese 2011 (Million-scale),
% Wong 2007 (Octant), Huffaker 2011 (Verfploeter), Geofeed RFCs.

\subsection{The IP Geolocation Toolkit for Operators}
\label{sec:background-toolkit}
% TODO: table comparing DNS-based, Geofeed, traceroute, and latency/CBG on pros/cons
% under operator constraints (scale, privacy, cost, freshness).
\textit{[Compares DNS-based, Geofeed, traceroute, and latency/CBG methods on scale, privacy, cost,
and freshness under operator constraints. Key point: latency methods double as validation for the
other three — operators use speed-of-Internet circles to sanity-check Geofeed/rDNS claims before
ingesting them as training labels.]}

\subsection{Constraint-Based Geolocation}
\label{sec:background-cbg}
% TODO: three-phase pipeline; notation; speed-of-light model; RTT → distance conversion.
\textit{[Three-phase pipeline: (1) LTD — per-VP latency-to-distance model maps RTT to a
geometric band; (2) MTL — intersect all VP constraints to form the feasible region R;
(3) CTR — select a point estimate from R. Inputs: RTT measurements + VP locations.
Outputs: estimated target lat/lon. Speed-of-Internet bound: $d \leq \frac{2}{3}c \cdot \text{RTT}/2$.]}

\subsection{The Four Foundational Variants}
\label{sec:background-variants}
% TODO: variant table (Vanilla, Million-scale, Octant, Spotter) with LTD / MTL / CTR columns.
% Implementation note: this benchmark uses planar_circle + geometric_centroid for
% Vanilla and Million-scale (not the original spherical-circle + BVM).

\begin{table}[h]
\caption{The four foundational CBG variants and their phase implementations.}
\label{tab:variants}
\begin{tabular}{llll}
\hline
\textbf{Variant} & \textbf{LTD (P1)} & \textbf{MTL (P2)} & \textbf{CTR (P3)} \\
\hline
Vanilla       & Low Envelope    & Planar circle             & Geometric centroid \\
Million-scale & Speed-of-Internet & Planar circle           & Geometric centroid \\
Octant        & Bounded Spline  & Planar annulus (weighted) & Monte Carlo medoid \\
Spotter       & Normal Distribution & Planar annulus (weighted) & Monte Carlo medoid \\
\hline
\end{tabular}
\end{table}

\textit{[Note: this benchmark uses \texttt{planar\_circle + geometric\_centroid} for Vanilla and
Million-scale. Speed-of-Internet (Million-scale) is the only calibration-free LTD variant — this
asymmetry drives the Anchored/Unanchored two-track taxonomy in \S\ref{sec:intro-challenges}.]}
